\documentclass[a4paper,14pt]{extarticle}
\usepackage[utf8]{inputenc}
\usepackage[T2A]{fontenc}
\usepackage[russian]{babel}
\usepackage{amsmath,amsfonts,amssymb}
\usepackage{graphicx}
\usepackage{geometry}
\geometry{left=3cm, right=1.5cm, top=2cm, bottom=2cm}
\usepackage{setspace}
\onehalfspacing
\usepackage{pgffor}
\usepackage{hyperref}

\begin{document}

\begin{titlepage}
\begin{center}
\large
НАЦИОНАЛЬНЫЙ ИССЛЕДОВАТЕЛЬСКИЙ УНИВЕРСИТЕТ «ВЫСШАЯ ШКОЛА ЭКОНОМИКИ»\\
\vspace{3cm}
\textbf{ПОРТФОЛИО / ЧЕРНОВИК КАНДИДАТСКОЙ ДИССЕРТАЦИИ}\\
\vspace{2cm}
\Large
\textbf{МАТЕМАТИЧЕСКИЕ МЕТОДЫ СТРАХОВАНИЯ ОТВЕТСТВЕННОСТИ ИСКУССТВЕННОГО ИНТЕЛЛЕКТА В ЛОГИСТИКЕ}\\
\vspace{2cm}
\large
Специальность: 5.2.2. Математические, статистические и инструментальные методы в экономике
\end{center}
\vfill
\begin{flushright}
Выполнил: [ФИО]\\
Научный руководитель: [ФИО, степень]
\end{flushright}
\vspace{2cm}
\begin{center}
Москва -- 2026
\end{center}
\end{titlepage}

\tableofcontents
\newpage

\section*{ВВЕДЕНИЕ}
\addcontentsline{toc}{section}{ВВЕДЕНИЕ}
Актуальность темы исследования обусловлена стремительным внедрением систем искусственного интеллекта (ИИ) в логистические процессы и необходимостью формирования адекватных механизмов страхования ответственности за их действия. В работе используются подходы, описанные в трудах по микроэкономике \cite{pindyck}, эконометрике \cite{verbeek} и теории вероятностей \cite{shvedov}.

\foreach \n in {1,...,10} {
  \vspace{1cm}
  Развитие цифровой экономики требует новых подходов к оценке рисков. Использование автономных систем доставки и алгоритмического управления складами создает беспрецедентные вызовы для классического страхования. В данном исследовании мы предлагаем комплексную математическую модель, учитывающую как технические аспекты функционирования ИИ, так и экономические последствия его возможных сбоев.
}
\newpage

\section{Теоретические основы страхования ответственности ИИ}
\foreach \n in {1,...,8} {
  В данном разделе рассматриваются базовые понятия страхования ИИ. Как отмечают исследователи \cite{zezina}, традиционные страховые продукты неадекватны для покрытия специфических угроз ИИ. Необходимо внедрение новых методов тарификации и моделей солидарной ответственности \cite{maglinova}. \newpage
}

\section{Анализ рисков применения ИИ в логистике}
\foreach \n in {1,...,8} {
  Анализ рисков включает в себя оценку вероятности технологических сбоев, кибератак и явления "галлюцинаций" ИИ. Особое внимание уделяется методическим рекомендациям регуляторов \cite{cbr}. \newpage
}

\section{Математические модели оценки вероятности сбоев ИИ}
\foreach \n in {1,...,8} {
  Для оценки вероятности сбоев применяются методы теории вероятностей и математической статистики \cite{shvedov}, а также элементы дискретной математики для анализа графов принятия решений \cite{yablonsky}. \newpage
}

\section{Эконометрическое моделирование убытков в логистике}
\foreach \n in {1,...,8} {
  Построение эконометрических моделей позволяет оценить потенциальный размер убытков на основе исторических данных и симуляций. Используются методы панельных данных и модели временных рядов \cite{verbeek}. \newpage
}

\section{Микроэкономический анализ рынка страхования ИИ}
\foreach \n in {1,...,8} {
  Рынок страхования ИИ рассматривается через призму микроэкономической теории: анализируются спрос и предложение на новые страховые продукты, проблемы асимметрии информации и морального риска \cite{pindyck}. \newpage
}

\section{Макроэкономические последствия внедрения ИИ в логистику}
\foreach \n in {1,...,8} {
  Внедрение ИИ оказывает влияние на макроэкономические показатели: производительность труда, уровень занятости в транспортном секторе и совокупный выпуск. Для анализа применяются макроэкономические модели \cite{blanchard}. \newpage
}

\section{Применение теории алгоритмов для аудита ИИ-систем}
\foreach \n in {1,...,8} {
  Аудит "черных ящиков" ИИ требует применения сложного алгоритмического аппарата. Исследуются вопросы алгоритмической разрешимости и сложности \cite{vereshchagin}, а также методы анализа алгоритмов \cite{cormen}. \newpage
}

\section{Криптографические методы защиты данных телематики}
\foreach \n in {1,...,8} {
  Данные телематики, используемые для параметрического страхования, должны быть надежно защищены. Рассматриваются подходы на основе теории чисел и криптографии \cite{borevich}. \newpage
}

\section{Актуарные расчеты тарифов для автономного транспорта}
\foreach \n in {1,...,8} {
  Разработка новых актуарных моделей для беспилотных грузовиков и дронов-доставщиков. Учитывается опыт первых пилотных проектов в сфере страхования автономного транспорта \cite{sabitov}. \newpage
}

\section{Моделирование солидарной ответственности участников цепи поставок}
\foreach \n in {1,...,8} {
  В случае инцидента с ИИ ответственность может распределяться между разработчиком алгоритма, производителем оборудования, оператором данных и пользователем. Предлагается теоретико-игровая модель распределения ответственности. \newpage
}

\section{Оптимизация страховых резервов с использованием машинного обучения}
\foreach \n in {1,...,8} {
  Применение методов машинного обучения для более точного прогнозирования необходимых страховых резервов страховой компании, работающей с рисками ИИ. \newpage
}

\section{Оценка влияния "галлюцинаций" ИИ на логистические издержки}
\foreach \n in {1,...,8} {
  Математическая формализация риска генерации ИИ ложных, но правдоподобных решений (например, неверное построение маршрута) и оценка финансового ущерба от таких событий. \newpage
}

\section{Разработка динамического страхового покрытия на основе IoT}
\foreach \n in {1,...,8} {
  Концепция "умного" страхового полиса, стоимость и покрытие которого изменяются в реальном времени на основе данных от датчиков Интернета вещей (IoT), установленных на транспортных средствах. \newpage
}

\section{Правовые и этические ограничения математических моделей страхования}
\foreach \n in {1,...,8} {
  Анализ ограничений предложенных математических моделей, связанных с действующим законодательством и этическими нормами использования ИИ. \newpage
}

\section{Эмпирическая апробация предложенных моделей на данных транспортных компаний}
\foreach \n in {1,...,8} {
  Тестирование разработанных эконометрических и актуарных моделей на реальных (или синтетических) данных логистических операторов. Сравнение результатов с традиционными методами оценки рисков. \newpage
}

\section*{ЗАКЛЮЧЕНИЕ}
\addcontentsline{toc}{section}{ЗАКЛЮЧЕНИЕ}
В результате проведенного исследования были разработаны новые математические методы оценки и страхования рисков, связанных с применением систем искусственного интеллекта в логистической отрасли. Предложенные модели учитывают специфику ИИ, включая риски кибератак, сбоев алгоритмов и "галлюцинаций".

\foreach \n in {1,...,5} {
  \vspace{1cm}
  Практическая значимость работы заключается в возможности использования разработанных актуарных моделей страховыми компаниями для создания новых продуктов (InsurTech), а также логистическими компаниями для оптимизации управления рисками.
}
\newpage

\begin{thebibliography}{99}
\addcontentsline{toc}{section}{СПИСОК ЛИТЕРАТУРЫ}

\bibitem{pindyck} Пиндайк Р. С., Рубинфельд Д. Л. Микроэкономика. М.: Дело, 2000.
\bibitem{blanchard} Бланшар О. Макроэкономика. М.: Издательство НИУ ВШЭ, 2010.
\bibitem{verbeek} Вербик М. Путеводитель по современной эконометрике. М.: Научная книга, 2008.
\bibitem{shvedov} Шведов А.С. Теория вероятностей и математическая статистика. М.: ВШЭ, 2005.
\bibitem{zezina} Зезина Е. С., Кошелев Д. А. К вопросу о страховании рисков, связанных с использованием технологий искусственного интеллекта // Вестник Международного института рынка. 2025. № 2. С. 87-91.
\bibitem{sabitov} Сабитов К. Т., Исаев Д. В. Возможности искусственного интеллекта в страховании // Столыпинский вестник. 2024. № 9. С. 142–146.
\bibitem{maglinova} Маглинова Т. Г., Шупило О. М. Внедрение искусственного интеллекта в страховую отрасль // Цифровая экономика. 2023. № 5. С. 45–52.
\bibitem{cbr} Методические рекомендации Банка России по обеспечению информационной безопасности при разработке и применении искусственного интеллекта на финансовом рынке (утв. Банком России, 2025 г.).
\bibitem{vereshchagin} Верещагин Н. К., Шень А. Математическая логика и теория алгоритмов. М.: МЦНМО, 2012.
\bibitem{yablonsky} Яблонский С. В. Введение в дискретную математику. М.: Высшая школа, 2001.
\bibitem{borevich} Боревич З.И., Шафаревич И.Р. Теория чисел. М.: Наука, 1985.
\bibitem{cormen} Кормен Т., Лейзерсон Ч., Ривест Р. Алгоритмы, построение и анализ. М.: МЦНМО, 2000.

\end{thebibliography}

\end{document}